% First we need to tell LaTeX what type of document we'll be writing
\documentclass{article} %lots of other options here, like amsart, among others

% Preamble:

%Here's where packages go. 
\usepackage{amsmath} %includes lots of mathematical symbols and content
\usepackage{lipsum} % For blind text
\usepackage{fullpage}
%\usepackage[margin=1in]{geometry}
%\usepackage{geometry}
%\geometry{
%  top=1in,
%  inner=1in,
%  outer=1in,
%  bottom=1in,
%}
\usepackage{enumerate}
\usepackage{array,tabularx,multirow} %Helpful for certain tables and table features
\usepackage{hyperref}%%Generally hyperref should be the last package loaded. 
%\hypersetup{
%    colorlinks=true,
%    urlcolor=blue,
%    linkcolor=black
%}


\usepackage[margin=1in]{geometry}
\usepackage{array}
\usepackage{multirow}
\title{Math 3680: Week 5}
\author{Arjan Singh Ellingson}
%\date{August 24th, 2026}
\date{\today}

\begin{document}

\maketitle
%\tableofcontents

\section{Equations}
I have taken this problem from a set of linear algebra problems I previously wrote for my textbook development position. However, previously, I did not have it in an equation environment. I may go back and standardize that!

\textbf{Example}: Find a linear combination of 
$\begin{bmatrix} 2 \\ 1 \\ -3 \end{bmatrix}$ and 
$\begin{bmatrix} 1 \\ -2 \\ 4 \end{bmatrix}$ 
that gives the vector 
$\begin{bmatrix} 17 \\ -4 \\ 2 \end{bmatrix}$.

\textbf{Solution}: We are looking for $a_1$ and $a_2$ such that:
\begin{equation*}
a_1 \begin{bmatrix} 2 \\ 1 \\ -3 \end{bmatrix}
+ 
a_2 \begin{bmatrix} 1 \\ -2 \\ 4 \end{bmatrix}
=
\begin{bmatrix} 17 \\ -4 \\ 2 \end{bmatrix}
\end{equation*}

Looking at each dimension separately, we get the system of equations:
\begin{align*}
2a_1 + a_2 &= 17 \\
a_1 - 2a_2 &= -4 \\
-3a_1 + 4a_2 &= 2
\end{align*}

If we can solve this system of equations, we will find $a_1$ and $a_2$. Let's 
multiply the first equation by 2 and add it to the second equation:
\begin{align*}
2\left(2a_1 + a_2\right) + \left(a_1 - 2a_2\right) &= 2(17) + (-4) \\
4a_1 + 2a_2 + a_1 - 2a_2 &= 34 - 4 \\
5a_1 &= 30 \\
a_1 &= 6
\end{align*}

Now we can take $a_1$ and substitute it back into any equation in our system 
to find $a_2$. Let's use the third equation:
\begin{align*}
-3(6) + 4a_2 &= 2 \\
-18 + 4a_2 &= 2 \\
4a_2 &= 20 \\
a_2 &= 5
\end{align*}



\section{Tables}

I had claude make a shortcut for me to create n-sized blocks in my table. A future improvement would be to reduce the amount of \texttt{\\cline} commands needed to make it a more readable and compact table.
% claude made a shortcut with two params: n 30-minute blocks and a label
% \block{#1}{#2}
% #1 = number of 30-min rows to span (e.g. 2 = 1 hour, 4 = 2 hours)
% #2 = label text

\newcommand{\block}[2]{\multirow{#1}{*}{\shortstack{#2}}}


% requires cline at the end of blocks: draws hlines only under provided columns.
\begin{table}[h]
\centering
\renewcommand{\arraystretch}{1.3}
\begin{tabular}{|>{\centering\arraybackslash}p{1.6cm}|
                 >{\centering\arraybackslash}p{2.2cm}|
                 >{\centering\arraybackslash}p{2.2cm}|
                 >{\centering\arraybackslash}p{2.2cm}|
                 >{\centering\arraybackslash}p{2.2cm}|
                 >{\centering\arraybackslash}p{2.2cm}|}
\hline
\textbf{Time} & \textbf{Mon} & \textbf{Tue} & \textbf{Wed} & \textbf{Thu} & \textbf{Fri} \\
\hline

8:00  & \block{4}{CSC 4880\\AI} & \block{4}{CSC 3300\\PL} & \block{4}{CSC 4880\\AI} & \block{4}{CSC 3300\\PL} & \block{4}{CSC 4880\\AI} \\ \cline{1-1}
8:30  & & & & & \\ \cline{1-1}
9:00  & & & & & \\ \cline{1-1}
9:30  & & & & & \\ \hline
10:00 & & & & & \\ \hline

10:30 & & \block{10}{Redwire \\(Work)} & & \block{10}{Redwire \\(Work)} & \\ \cline{1-2}\cline{4-4}\cline{6-6}
11:00 & \block{2}{MATH 3680\\\LaTeX} & & & & \\ \cline{1-1}\cline{4-4}\cline{6-6}
11:30 & & & & & \\ \cline{1-2}\cline{4-4}\cline{6-6}

12:00 & & & & & \\ \cline{1-2}\cline{4-4}\cline{6-6}
12:30 & & & & & \\ \cline{1-2}\cline{4-4}\cline{6-6}

1:00  & & & & & \\ \cline{1-2}\cline{4-4}\cline{6-6}
1:30  & \block{6}{CSC 4160\\SE Capstone} & & \block{6}{CSC 4160\\SE Capstone} & & \\ \cline{1-1}\cline{6-6}
2:00  & & & & & \\ \cline{1-1}\cline{6-6}
2:30  & & & & & \\ \cline{1-1}\cline{6-6}
3:00  & & & & & \\ \cline{1-1}\cline{3-3}\cline{5-6}
3:30  & & & & & \\ \cline{1-1}\cline{3-3}\cline{5-6}
4:00  & & & & \block{2}{Front Porch \\ Volunteering!!} & \\ \cline{1-4}\cline{6-6}

4:30  & & & & & \\ \hline

5:00  & \block{4}{CSC 3665\\Databases} & & \block{4}{CSC 3665\\Databases} & & \block{4}{CSC 3665\\Databases} \\ \cline{1-2}\cline{3-4}\cline{5-6}
5:30  & & & & & \\ \cline{1-1}\cline{3-3}\cline{5-5}
6:00  & & & & & \\ \cline{1-1}\cline{3-3}\cline{5-5}
6:30  & & & & & \\ \hline

7:00  & & & & & \\ \hline
7:30  & & & & & \\ \hline

8:00  & & & & & \\ \hline

\end{tabular}
\caption{Weekly Schedule}
\end{table}



\end{document}
